\scanpage{72}
% loc: 1
et on a un \deleted{bijection} diagramme commutatif

% loc: 2–5, upper diagram
\diagram{\begin{tikzcd}[ampersand replacement=\&,column sep=large,row sep=huge]
 {\left[E^*/k^*\underset{\substack{\vert\\\text{si }\operatorname{Pic}(k)=0\\\text{sinon }\hookrightarrow}}{\simeq}\mathbb P(E)(k)\right]}
 \arrow[r,"{\substack{\text{induit}\\\deleted{\text{\unsure{fonctoriel|fonctif}}}\,\text{iso de}\\\text{schémas}\,\deleted{\unsure{\sim}}\\\sim}}","{\text{bien connu}}"']
 \& C^*/k^* \& {}\\
 E^*\arrow[u]\arrow[r,"{\varphi^e}"]
 \& C^*(e)\arrow[u]\arrow[r,phantom,"{\subset}"]
 \& C^\times\arrow[ul,dashed]
\end{tikzcd}}

% loc: 6–7
On a $\varphi^e(x)=\varphi^e(y)$ \unsure{ses|sés|ssi} $y=\lambda x$ avec $\lambda\in\mu_2(k)$ (i.e. $\lambda\in k^\times$, $\lambda^2=1$).

% loc: 8–9
Supposons maintenant
\[
 \{\pm1\}\xrightarrow{\overset{\text{iso}}{\sim}}\mu_2(k).
\]

% loc: 10–16
[Si 2 inv. Cela signifie que $\operatorname{Spec}k$ est 0-connexe. Si $\operatorname{Spec}k$ a un nb fini \inserted{$n$} de comp. conn. \deleted{\unsure{soient}}, cela signifie que \deleted{toutes} le nb de comp. connexes où $2$ est nul est $n-1$ -- i.e.$^*$ $k=k'\times k''$ \deleted{\unsure{avec}} \unsure{anneau produit|anneau premier}, avec $2$ nul dans $k'$, \uline{non} nul dans $k''$, et $\operatorname{Spec}(k'')$ est \deleted{\unsure{co}} connexe.]

% loc: left beside 13–16
% slanted
\margin{le i.e. est sans doute vrai sans hyp. sur le nb de comp. connex.}

% loc: 17–20, bottom
Alors l'hom induit
\[
 \psi^e:\quad E^*/{\pm1}\xrightarrow{\sim}C^*(e)
\]
est \uline{bijectif}. Cela nous permet de transporter sur $C^*(e)$ la structure
% archive sheet number: 71

